% Files should be in encoding corresponding to the setting of \usepackage[...]{inputenc}

\documentclass[%      Basic settings of the document
%  draft,    				  % Turn on draft compilation
  14pt,       				% 12 points default font size
	c,                  % The slides should always start at their top (not vertically centered)
	aspectratio=1610,   % Ratio of width and height is 16:10 (corresponding to lecture rooms at Technická 12),
	                    % other options are 43, 149, 169, 54, 32.
	unicode,						% Bookmarks and metainformation in the compiled PDF will be unicode
]{beamer}				    	% Document type 'beamer' is the most often used one for slides
%\usepackage{etex}

\usepackage{hyperref}
\usepackage[utf8]		  %	The input files are in UTF-8 (unicode)
	{inputenc}					% Package for setting input encoding
	
\usepackage{graphicx} % Package for inclusion of external graphics

\usepackage[
	nohyperlinks				% No hyperlinks to the list of acronyms will be created
]{acronym}						% Package 'acronym' for automated handling of abbreviations and symbols
											% Neccessary for the 'acronym' environment of 'thesis' to work

\usepackage{cmap} 		% This package allows the final PDF
											% is fully "searchable" and "copyable"

%\usepackage{amsmath} %for typesetting advanced mathematics

\usepackage{booktabs} % Package allowing the use of commands \toprule, \midrule, \bottomrule
                      % when creating tables


%%%%%%%%%%%%%%%%%%%%%%%%%%%%%%%%%%%%%%%%%%%%%%%%%%%%%%%%%%%%%%%%%
%%%%%%%%                   Definitions                 %%%%%%%%%%
%%%%%%%%%%%%%%%%%%%%%%%%%%%%%%%%%%%%%%%%%%%%%%%%%%%%%%%%%%%%%%%%%

% Variable fields such as your name, title of the thesis atc. are set in this file.
% The file is SHARED between the thesis text and the presentation --- no need to set anything twice.

\usepackage[
  english,			% English study program
%%% Choose one of the thesis types
	semestral,		%	semestral project (abstracts, declaration and acknowledgements are excluded) (default)
  %bachelor,			%	Bachelor's thesis
  %master,			% Master's thesis
  %treatise,			% Treatise on doctoral thesis
  %doctoral,				% Doctoral thesis
%%% Choose one of the options:
%  left,				% Equations and captions will be aligned to the left
  center,			% Equations and captions will be centered (default)
]{thesis}   % Package for typesetting theses


%%% First and last name of the thesis author, following the scheme
% [titles in front]{FirstName}{LastName}[titles after]
% If the author does not have a title in front/after, just delete the entire string '[...]'
\author[Bc.]{Samuel}{Kopecký}

%%% Brno University of Technology identification number of author (BUT ID)
\butid{007}

%%% First and last name of the advisor
% [titles in front]{FirstName}{LastName}[titles after]
% If the author does not have a title in front/after, just delete the entire string '[...]'
% [titles in front]{FirstName}{LastName}[titles after]
\advisor[prof.\ Ing.]{FirstName}{LastName}[CSc.]

%%% First and last name of the opponent
% [titles in front]{FirstName}{LastName}[titles after]
% Makes use only in the presentation for defense;
% in case you do not want the opponent to be shown on the title slide, just comment out the command
% The opponent is not shown in case of the semestral thesis (no opponent exists actually at that time)
% In the case of a doctoral thesis, two to three oponents are typically involved. In such a case, if you want to have them on the title slide, please go to the definition of "VUT title page" in the thesis.sty file, uncomment and adapt their names.
\opponent[doc.\ Mgr.]{FirstName}{LastName}[Ph.D.]

%%% Thesis title
%  In the case of a very long thesis title, it may happen that it does not fit
%  into the slot in the footbar of the slides. You may use the command
%  \def\insertshorttitle{Shortened th.\ title}
%  where the shortened version of the title appears as the parameter.
%  If you do not want to shorten the title, you will have to redefine how the slide footbar
%  is generated, see: https://bit.ly/3EJTp5A
\title{Post-Quantum One-Time Linkable Ring Signature}

%%% Study program/specialization
\specialization{Teleinformatics} %Teleinformatika

%%% Department
%\department{Department of Control and Instrumentation}  %Ústav automatizace a měřicí techniky
%\department{Department of Biomedical Engineering}       %Ústav biomedicínského inženýrství
%\department{Department of Electrical Power Engineering} %Ústav elektroenergetiky
%\department{Department of Electrical and Electronic Technology}   %Ústav elektrotechnologie
%\department{Department of Physics}                      %Ústav fyziky
%\department{Department of Foreign Languages}            %Ústav jazyků
%\department{Department of Mathematics}                  %Ústav matematiky
%\department{Department of Microelectronics}             %Ústav mikroelektroniky
%\department{Department of Radio Electronics}            %Ústav radioelektroniky
%\department{Department of Theoretical and Experimental Electrical Engineering}  %Ústav teoretické a experimentální elektrotechniky
\department{Department of Telecommunications}           %Ústav telekomunikací
%\department{Department of Power Electrical and Electronic Engineering}   %Ústav výkonové elektrotechniky a elektroniky

%%% Faculty
%\faculty{Faculty of Architecture}   %Fakulta architektury
\faculty{Faculty of Electrical Engineering and~Communication}   %Fakulta elektrotechniky a~komunikačních technologií
%\faculty{Faculty of Chemistry}   %Fakulta chemická
%\faculty{Faculty of Information Technology}   %Fakulta informačních technologií
%\faculty{Faculty of Business and Management}   $Fakulta podnikatelská
%\faculty{Faculty of Civil Engineering}   %Fakulta stavební
%\faculty{Faculty of Mechanical Engineering}   %Fakulta strojního inženýrství
%\faculty{Faculty of Fine Arts}   %Fakulta výtvarných umění
%
%Logotype selection (in square brackets short logo, in curly brackets full logo):
\facultylogo[logo/FEEC_abbreviation_color_PANTONE_EN]{logo/UTKO_color_PANTONE_EN}


%%% Graduate year (typically the calendar year of the defense)
\graduateyear{2023}
%%% Academic year (typically the year of solution of the thesis in the format n/n+1)
\academicyear{2022/23}
% Date of the defense (makes use only in the presentation slides)
\date{13 January 2023} 

%%% Place of the defense
\city{Brno}

%%% Abstract
\abstract{%
Abstract in English.
}

%%% Keywords
\keywrds{%
Keywords in English
}

%%% Thanks and acknowledgement
\acknowledgement{%
I would like to thank the advisor of my thesis, Ing.~XXX YYY, Ph.D.\ for his/her valuable comments etc.
}%      % please go into that file and fill information about you, your topic etc.


%%%%%%%%%%%%%%%%%%%%%%%%%%%%%%%%%%%%%%%%%%%%%%%%%%%%%%%%%%%%%%%%%%%%%%%%
%%%%%   Setting values that appear in the Properties of the PDF  %%%%%%%
%%%%%%%%%%%%%%%%%%%%%%%%%%%%%%%%%%%%%%%%%%%%%%%%%%%%%%%%%%%%%%%%%%%%%%%%
%% If 'hyperref' package is loaded then command '\pdfsettings' can be simply used
\pdfsettings
%  However, you can set the individual fields by hand:
%\hypersetup{
%  pdftitle={Title of the thesis},    	% 'Document Title' field
%  pdfauthor={Author name},   	        % 'Author' field
%  pdfsubject={Type of the document}, 	% 'Subject' field
%  pdfkeywords={Keywords}           	  % 'Keywords' field
%}
\hypersetup{pdfpagemode=FullScreen}  % when opening the file with Adobe Reader,
                                     % automatically open it in fullscreen mode
%%%%%%%%%%%%%%%%%%%%%%%%%%%%%%%%%%%%%%%%%%%%%%%%%%%%%%%%%%%%%%%%%%%%%%%

\usetheme{VUT} 				% colours and layout corresponding to the BUT style
% alternatively, other styles can be used, such as
%\usetheme{Darmstadt} \usecolortheme{default2}
% but with no guarantees
\logoheader					% inserts small BUT FEEC logotype in the header; leave this in use




\begin{document}

% if you do not use the following command, the lower right corner of each slide will contain active navigation symbols
\disablenavigationsymbols

% the title slide will be created without header, footer etc.
\maketitle

\begin{frame}
  \frametitle{Table of contents}
  \large{
    \begin{itemize}
      \item Introduction
      \item Article's contribution
      \item Proposed solution
      \item Results
      \item Conclusion
    \end{itemize}
  }
\end{frame}

\begin{frame}
  \frametitle{Ring signatures}
  \large{
    \begin{columns}[T]
      \begin{column}{0.8\textwidth}
        \begin{itemize}
          \item Asymmetric cryptography
          \item Derived from group signatures
          \item Sign a message on behalf of a group of signers
          \item Preserves the anonymity of the signer
          \item Used in \textbf{cryptocurrencies}, \textbf{e-voting}
        \end{itemize}
      \end{column}
      \begin{column}{0.2\textwidth}
        \begin{figure}[htbp]
          \centering
          \includegraphics[width=\textwidth]{images/monero.png}
        \end{figure}
      \end{column}
    \end{columns}
  }
\end{frame}

\begin{frame} 
  \frametitle{Linkability of ring signatures}
  \large{
    \begin{itemize}
      \item Known as Linkable Ring Signatures
      \item A signer generates a label (tag)
      \item A label is generated from the private key
      \item Detect whether two signatures were generated by the \textbf{same signer}
      \item Used in Monero to prevent double spending
    \end{itemize}
  }
\end{frame}

\begin{frame} 
  \frametitle{Post-quantum cryptography}
  \large{
    \begin{itemize}
      \item Currently used algorithms can be broken using quantum computers
      \item Applies mostly to asymmetric cryptography
      \item Best quantum computer has \textbf{127 qubits} (IBM)[1]
      \item \textbf{Tens of millions} bits required to efficiently break asymmetric cryptography
    \end{itemize}
  }
\end{frame}

\begin{frame}
  \frametitle{Proposed solution}
  \large{
    \begin{itemize}
      \item A need to create post-quantum ring signatures
      \item Authors introduce \textbf{Lattice-Based one-time Linkable Ring signature} (L2RS) scheme
      \item Cryptocurrency privacy-preserving protocol called \textbf{Lattice RingCT v1.0}
      \item Monero currently uses Ring CT
    \end{itemize}
  }
\end{frame}

\begin{frame}
  \frametitle{About L2RS scheme}
  \large{
    \begin{itemize}
      \item Based on Bimodal Lattice Signature Scheme (BLISS)
      \item Relies on the \textbf{Ring-SIS} (Short Integer Solution) problem
      \item It is secure in terms of \textbf{unforgeabillity}, \textbf{linkability}, \textbf{non-slanderability} 
      \item It also achieves \textbf{unconditional anonymity}
    \end{itemize}
  }
\end{frame}

\begin{frame}
  \frametitle{L2RS foundations}
  \large{
    \begin{itemize}
      \item The main structure is a \textbf{ring} $R_q = \mathbb{Z}_q[X]/f(x)$
      \item The number $q$ is a prime number where $q=1\,\mathrm{mod}\,2n$ and $n=512$
      \item \textbf{Private key} consists of $m-1$ rings (polynomials) where $f(x) = x^n - 1$ and coefficients modulo $q$ 
      \item \textbf{The public key} consists of only one ring
      \item \textbf{2 Public parameters} are the same size as the private key
    \end{itemize}
  }
\end{frame}

\begin{frame}
  \frametitle{Rings}
  \large{
    An example of a private key \textbf{S} for $m=6$, $n=512$
    \begin{itemize}
      \item $s_1 = s_{1,0} + s_{1,1}x + s_{1,2}x^2 + \dots + s_{1,510}x^{510} + s_{1,511}x^{511}$
      \item $s_2 = s_{2,0} + s_{2,1}x + s_{2,2}x^2 + \dots + s_{2,510}x^{510} + s_{2,511}x^{511}$
      \item $s_3 = s_{3,0} + s_{3,1}x + s_{3,2}x^2 + \dots + s_{3,510}x^{510} + s_{3,511}x^{511}$
      \item $s_4 = s_{4,0} + s_{4,1}x + s_{4,2}x^2 + \dots + s_{4,510}x^{510} + s_{4,511}x^{511}$
      \item $s_5 = s_{5,0} + s_{5,1}x + s_{5,2}x^2 + \dots + s_{5,510}x^{510} + s_{5,511}x^{511}$
    \end{itemize}
  }
\end{frame}

\begin{frame}
  \frametitle{L2RS signing/verifying}
  \large{
    \begin{itemize}
      \item $w$ number of participants
      \item All the public \textbf{keys are colleted} including the signers
      \item The signatures consists of a hash $c_1$ and rings $t_1,\dots,t_w$
      \item It is created using all the \textbf{public keys} the \textbf{signers private key} and \textbf{the public parameters}
      \item The signature can verified using the \textbf{public keys} and \textbf{public parameters} while keeping the signers identity
    \end{itemize}
  }
\end{frame}

\begin{frame}
  \frametitle{Results}
  \large{
    \begin{itemize}
      \item Many of the security claims are well described in the papers appendixes
      \item This scheme was used to build an extension of the original Ring CT protocol used by Monero cryptocurrency
      \item There is no reference implementation by the authors so there are no to performance results
    \end{itemize}
  }
\end{frame}

\begin{frame}
  \frametitle{Conclusion}
  \large{
    \begin{itemize}
      \item The authors successfully managed to create a secure post quantum ring signature scheme
      \item The L2RS scheme can be used in many areas not only cryptocurrencies as showcased in this paper
    \end{itemize}
  }
\end{frame}

\begin{frame}
  \frametitle{Bibliography}
  \small{
    [1] IBM Unveils Breakthrough 127-Qubit Quantum Processor. IBM. \textit{IBM Newsroom} [online]. 2021 [cit. 2022-10-26]. Available from: \url{https://shorturl.at/vDJT4}
  }
\end{frame}

\begin{frame}[c]
  \frametitle{\mbox{ }}
  \begin{center}
    {\Huge Thank you for your attention!}
  \end{center}
\end{frame}

\end{document}
